\documentclass[10pt]{article}
\usepackage[margin=1in]{geometry}
\usepackage[T1]{fontenc}
\usepackage{lmodern,microtype,amsmath,amssymb,amsthm}
\theoremstyle{definition}
\newtheorem{assumption}{Assumption}
\begin{document}
\section*{Sampled Smooth Convex Optimization}
We minimize $f(x)=\mathbb{E}_{\xi\sim\mu}[\ell(\xi,x)]$ on a separable real Hilbert space $E$, where $\mu$ is a probability law on a measurable sample space $\Xi$. Each sampled loss $\ell(\xi,\cdot)$ is convex with $L$-Lipschitz gradient, $L>0$, and losses are integrable. A given $x^0\in E$ is within $R\geq0$ of a global minimizer $x^\star$, where $\mathbb{E}_{\xi\sim\mu}\|\nabla_x\ell(\xi,x^\star)\|^2\leq\sigma^2$ for $\sigma\geq0$. Each call returns a sampled gradient using a fresh independent sample. For $\varepsilon>0$, success requires a measurable output $\bar{x}$ with integrable gap and $\mathbb{E}[f(\bar{x})-f(x^\star)]\leq\varepsilon$, whose law equals that of a linearly decaying stochastic-gradient last iterate.

The optimization problem is
\[
\min_{x\in E} f(x),
\qquad
f(x)=\int_{\Xi}\ell(\xi,x)\,\mu(d\xi).
\]
The norm on $E$ is induced by its inner product. The public inputs are $E$, $\Xi$, $L$, $R$, $\sigma$, $\varepsilon$, and $x^0$, with the parameter ranges stated above. The probability law $\mu$, sampled losses $\ell$, and sampled-gradient map $g:\Xi\times E\to E$ belong to the unknown instance. They satisfy the following conditions.

\begin{assumption}[Sampled convexity and smoothness]
For every $\xi\in\Xi$, the function $\ell(\xi,\cdot)$ is differentiable everywhere. For every $x,y\in E$ and $t\in[0,1]$,
\begin{align*}
g(\xi,x)&=\nabla_x\ell(\xi,x),\\
\ell(\xi,(1-t)x+ty)&\leq(1-t)\ell(\xi,x)+t\ell(\xi,y),\\
\|g(\xi,x)-g(\xi,y)\|&\leq L\|x-y\|.
\end{align*}
\end{assumption}

\begin{assumption}[Measurability and loss integrability]
The map $(\xi,x)\mapsto g(\xi,x)$ is jointly measurable. For every $x\in E$, the function $\xi\mapsto\ell(\xi,x)$ is integrable under $\mu$.
\end{assumption}

\begin{assumption}[Attainment and gradient moment]
There exists $x^\star\in E$ such that
\begin{align*}
f(x^\star)&\leq f(x)\qquad\text{for every }x\in E,\\
\|x^0-x^\star\|&\leq R.
\end{align*}
At this same minimizer, $\xi\mapsto\|g(\xi,x^\star)\|^2$ is integrable and
\[
\int_{\Xi}\|g(\xi,x^\star)\|^2\,\mu(d\xi)\leq\sigma^2.
\]
\end{assumption}

\begin{assumption}[Independent sampled-gradient access]
At call index $k\geq0$, a query $x\in E$ returns $g(\xi_k,x)\in E$, where $(\xi_k)_{k\geq0}$ has law $\mu^{\otimes\mathbb{N}}$. Each query costs one oracle call. Queries may be chosen adaptively anywhere in $E$, and any internal algorithm randomness is independent of the sample sequence.
\end{assumption}

An externally specified budget $K\in\mathbb{N}$ may depend on the full public parameter tuple. The algorithm uses at most $K$ oracle calls and returns a point $\bar{x}\in E$. Besides the accuracy requirement, success imposes a specific output-distribution condition. For a deterministic integer $T\geq0$, define
\[
\eta_T=\frac{R}{8LR+\sigma\sqrt{T}},
\]
with the quotient taken as zero when its denominator is zero. For independent samples $(\xi_k)_{k\geq0}$ with common law $\mu$, define the iterates by
\begin{align*}
x_T^0&=x^0,\\
x_T^{k+1}&=x_T^k-\eta_T\frac{T-k}{T}\,g(\xi_k,x_T^k),
\qquad 0\leq k<T.
\end{align*}
For $T=0$, the last iterate is $x_0^0=x^0$. Success requires the existence of such a $T$ for which the law of $\bar{x}$ equals the law of $x_T^T$; the distributional condition does not identify $T$ with the call budget $K$.

Writing $\mathbb{P}$ for the joint probability law of the algorithm and oracle randomness, the output must be measurable. There must also exist a global minimizer $x^\star$ such that the objective gap is integrable under $\mathbb{P}$ and
\[
\int \bigl(f(\bar{x}(\omega))-f(x^\star)\bigr)\,\mathbb{P}(d\omega)
\leq\varepsilon.
\]
Since all global minimizers have the same objective value, the minimizer witnessing the initial distance and gradient-moment conditions may be used in this accuracy criterion.
\end{document}
